\documentclass[11pt,a4paper]{article}
\usepackage[utf-8]{inputenc}
\usepackage[margin=1in]{geometry}
\usepackage{amsmath}
\usepackage{amssymb}
\usepackage{graphicx}
\usepackage{booktabs}
\usepackage{hyperref}
\usepackage{fancyhdr}
\usepackage{lastpage}

% Header and Footer
\pagestyle{fancy}
\fancyhf{}
\rhead{Mahdi Zeineldin}
\lhead{ResNet-9 Small Dataset Learning}
\cfoot{Page \thepage\ of \pageref{LastPage}}

% Title
\title{\Large \textbf{Improving ResNet-9 Generalization Trained on Small Datasets}}
\author{Mahdi Zeineldin \\ zeinaldinmahdi@gmail.com \\ Lebanese American University}
\date{August 20, 2026}

\begin{document}

\maketitle

\section*{Abstract}

This project successfully reproduces and extends the ICLR 2023 paper ``Improving ResNet-9 Generalization Trained on Small Datasets'' by combining Sharpness Aware Minimization (SAM), Label Smoothing, Gradient Centralization, and Input Patch Whitening to improve ResNet-9 performance on CIFAR-10 with only 10\% of training data (5,000 images). The baseline model achieves 74.45\% test accuracy, while the optimized model reaches 86.09\%, representing an 11.64\% improvement that closely matches the paper's target of $\sim$88\%. As an extension, robustness testing reveals that SAM-optimized models demonstrate superior resilience to various perturbations, with average robustness improvements of 6.7\% across multiple noise types. The results validate the effectiveness of combining multiple regularization and optimization techniques for improved generalization on small datasets.

\section{Introduction \& Problem Statement}

Deep learning models often underperform when training data is limited, presenting a critical challenge in real-world applications where collecting large labeled datasets is prohibitively expensive. This problem is particularly acute in medical imaging, robotics, and domain-specific applications with limited labeled data availability.

The paper we reproduce addresses this gap by proposing a combination of four complementary optimization techniques. Our core research questions are: (1) How effective are SAM, label smoothing, gradient centralization, and patch whitening when combined? (2) Do these techniques improve model robustness beyond accuracy? The project goal is to reproduce the paper's results (target: $\sim$88\% accuracy on CIFAR-10 with 10\% training data) and extend the work by evaluating robustness to various perturbations.

\section{Methodology}

\subsection{Architecture: ResNet-9}
We implement a lightweight ResNet-9 architecture designed for small datasets. The model consists of a preparatory layer (Conv 3$\rightarrow$64), three residual layers (Layer 1: 64$\rightarrow$64, Layer 2: 64$\rightarrow$128 with stride 2, Layer 3: 128$\rightarrow$256 with stride 2), global average pooling, and a final fully-connected layer. Total parameters: $\sim$2.5M.

\subsection{Optimization Techniques}

\textbf{Sharpness Aware Minimization (SAM):} Seeks flat loss minima by performing two-step optimization: first computing loss and gradients at current weights, then computing loss at perturbed weights (ρ=0.05). This addresses the observation that sharp minima often lead to poor generalization.

\textbf{Label Smoothing:} Converts hard one-hot targets to soft targets (correct class: 0.9, wrong classes: 0.1/9). This regularization reduces overfitting by preventing the model from becoming overconfident.

\textbf{Gradient Centralization:} Removes gradient correlation by subtracting layer-wise mean gradients, smoothing the loss landscape and improving optimization dynamics.

\textbf{Input Patch Whitening:} Normalizes 4$\times$4 image patches using (x - mean)/std transformation, reducing input redundancy and improving feature learning.

\subsection{Dataset \& Training Configuration}
Dataset: CIFAR-10 with 10\% subset (5,000 training images, 10,000 test images). Augmentation includes random crop and horizontal flip. Hyperparameters: learning rate 0.1, momentum 0.9, weight decay 5e-4, batch size 128, 50 epochs.

\section{Implementation Details \& Results}

\subsection{Experimental Setup}
Two configurations were trained: (1) Baseline using SGD with cross-entropy loss, and (2) Optimized using SAM with label smoothing, gradient centralization, and patch whitening.

\subsection{Quantitative Results}

\begin{table}[h]
\centering
\begin{tabular}{lccc}
\toprule
\textbf{Configuration} & \textbf{Accuracy} & \textbf{vs. Paper} & \textbf{Time (min)} \\
\midrule
Baseline (SGD + CE) & 74.45\% & Target $\sim$76\% & 8.5 \\
Optimized (SAM+LS+GC+PW) & 86.09\% & Target $\sim$88\% & 9.2 \\
\textbf{Improvement} & \textbf{+11.64\%} & \textbf{Match: +10-12\%} & \textbf{+8\%} \\
\bottomrule
\end{tabular}
\end{table}

\subsection{Ablation Study}
Individual technique contributions: SAM (+4.7\%), Gradient Centralization (+2.4\%), Label Smoothing (+1.8\%), Patch Whitening (+0.9\%). The combined effect (11.64\%) exceeds the sum, demonstrating synergistic benefits.

\subsection{Extension: Robustness Analysis}
Testing model resilience to perturbations reveals consistent improvements for the optimized model:
\begin{itemize}
    \item Gaussian noise (σ=0.20): Baseline 48.19\% → Optimized 58.73\% (+10.54\%)
    \item Salt-and-pepper noise: Baseline 71.23\% → Optimized 76.89\% (+5.66\%)
    \item Brightness variation: Baseline 74.15\% → Optimized 79.34\% (+5.19\%)
\end{itemize}
Average robustness improvement: 6.7\% across all perturbations.

\section{Discussion \& Analysis}

Our results closely match the paper's findings, with optimized accuracy of 86.09\% versus the target of $\sim$88\% (1.91\% difference). SAM emerges as the dominant contributor, suggesting that finding flat loss minima is crucial for small-dataset generalization. The synergistic effects of combining multiple techniques indicate they address complementary aspects of the learning problem.

The robustness extension reveals that optimization for generalization also improves feature robustness to perturbations—a valuable finding suggesting that regularization techniques provide dual benefits.

\textbf{Limitations:} Evaluation is limited to CIFAR-10; cross-dataset validation would strengthen claims. Hyperparameter sensitivity (particularly ρ for SAM and smoothing coefficient) merits investigation. The techniques are optimized for 10\% data; behavior on larger datasets requires study.

\textbf{Future Work:} (1) Transfer learning experiments with ImageNet pre-training, (2) Application to medical imaging and robotics datasets, (3) Theoretical analysis of synergistic interactions, (4) Systematic hyperparameter optimization.

\section{Reflection on Learnings}

\textbf{Rewarding Aspects:} Successfully reproducing published results within 2\% of the target demonstrated the importance of careful implementation and attention to detail. The robustness extension added novel research value by revealing that optimized models are inherently more robust—an insight not present in the original paper.

\textbf{Key Challenges \& Solutions:}
\begin{itemize}
    \item \textit{SSL Certificate Issues:} Disabled SSL verification for dataset download (common on some systems)
    \item \textit{Hyperparameter Tuning:} Validated against paper values rather than grid searching to maintain reproducibility
    \item \textit{Robustness Testing Complexity:} Separated into extension phase to keep main pipeline focused
\end{itemize}

\textbf{Skill Development:} This project deepened expertise in deep learning optimization (SAM, gradient manipulation), experimental design (ablation studies), research reproducibility, and scientific communication. The modular code structure enables future extensions in transfer learning and domain adaptation.

\section*{Conclusion}

We successfully reproduced the ICLR 2023 paper's core results and extended the work with robustness analysis. The 86.09\% accuracy (11.64\% improvement over baseline) validates the effectiveness of combining SAM, label smoothing, gradient centralization, and patch whitening for small-dataset learning. The extension demonstrates that these techniques provide additional benefits beyond accuracy—improved model robustness to perturbations. The clean, modular implementation and comprehensive documentation enable reproducibility and future research extensions.

\end{document}
